%!TEX TS-program = xelatex
%!TEX encoding = UTF-8 Unicode

\documentclass[english,a4paper,11pt]{article}

\nonstopmode

%\usepackage[T1]{fontenc}
%\usepackage[utf8]{inputenc}
\usepackage[margin=3cm,a4paper]{geometry}

\usepackage{graphicx}
\usepackage{amssymb}
\usepackage[norsk]{babel}
\usepackage{color}
\usepackage{xcolor}
\usepackage{booktabs}
\usepackage{array}
\usepackage{wrapfig}
\usepackage{calc}
\usepackage{hyperref}
\usepackage{xspace}
\usepackage{parskip}


\definecolor{xalblue}{rgb}{0.1992,0.1992,0.5977}
\definecolor{xaldark}{rgb}{0.1992,0.1992,0}


\usepackage{fontspec,xltxtra,xunicode,relsize}
\defaultfontfeatures{Mapping=tex-text}
\setmainfont[Mapping=tex-text]{Verlag Light}
%\setmainfont[Mapping=tex-text]{Sina-Regular}
\setsansfont[Scale=MatchLowercase,Mapping=tex-text]{Helvetica Neue}
\setmonofont[Scale=MatchLowercase]{Monaco}

\definecolor{lightgray}{gray}{0.8}
\newcolumntype{L}{>{\raggedright}p{0.14\textwidth}}
\newcolumntype{R}{p{0.76\textwidth}}
\newcommand\VRule{\color{lightgray}\vrule width 0.5pt}


\usepackage{titlesec}
\newcommand\secformat[1]{\normalfont\Large\bfseries\color{xalblue}#1}
\titleformat{name=\section}[block]{}{}{}{\secformat}

\begin{document}

\vspace*{-3cm}
\hspace*{12.5cm}\includegraphics[width=4cm]{logo.pdf}

\vspace*{-1.65cm}

\noindent
{\LARGE \bfseries \color{xaldark}Curriculum Vitae for Åsmund Ødegård}

\section*{Personal information}

\begin{minipage}[t]{0.4\textwidth}
\begin{tabbing}Nationality---\=olaf\kill 
Address: \>Sneisveien 26, \\
\>1407 Vinterbro \\
Born:\>23.08.1972 \\
\end{tabbing}
\end{minipage}
$\quad\quad$
\begin{minipage}[t]{0.3\textwidth}
\begin{tabbing}Nationality---\=olaf\kill 
E-mail: \>asmund@xal.no\\
Phone: \>+47 900 69 915\\
Nationality:\>Norwegian\\
\end{tabbing}
\end{minipage}

\section*{Summary}

I am a senior software developer and project leader. Most of my professional
career has been with Simula Research Laboratory, 
recently ranked as Norway's best ICT research institution. Through various
positions in the company I have obtained a broad set of technical and management
skills. My main areas of technical expertise spans from mathematical 
software development and scripting, to administration of advanced ICT
solutions and configuration management. Both through education and
practical hands-on experience I have developed strong analytical skills within
my fields of expertise. 

As manager and leader of several projects with internal and external
participants, I have learned the importance of understanding the strengths and
weaknesses of the participants in order to deliver on time. Communication within
the team is always key to ensure that everyone have ownership in the project and
stays motivated. This must always be balanced with a firm leadership.

\section*{Technical Skills}

\begin{tabular}{LR}
Languages & Python, C, C++, PHP, Java, Javascript, Bash, Matlab, Perl,
LaTeX, HTML, Sql, CSS, Scala, Common-Lisp, Clojure \\

Frameworks & jQuery, Web.py, Turbogears, Numpy, Pandas, SciPy, MPI, BSP,
Diffpack, Bootstrap, Plone, buildout, Puppet \\

Tools & Linux, OS X, VirtualBox, Vagrant, MySql, Apache, Nginx, Subversion, CVS,
Git, Mercurial, Make, SCons, Autoconf \\
\end{tabular}

\section*{Education}
\begin{tabular}{LR}
2006 & PhD (Dr. Scient) in Scientific Computing, 
University of Oslo, Norway. Thesis' title ``Applications of high level software
for parallel solution of Partial Differential Equations'' \\
1995 & Cand. Scient in Applied and Industrial Mathematics, University of
Oslo, Norway.
\end{tabular}


\section*{Professional Experience}

\begin{tabular}{LR}
2011 -- 2013 &  Director of External Research Funding at Simula Research
   Laboratory \\
2010 -- 2011 & Director of Administration at Simula Research Laboratory \\
2008 -- 2010 & Assistant Director of Administration at Simula Research Laboratory \\
2005 -- 2007 & Research Scientist at Simula Research Laboratory \\ 
2004 -- 2009 & IT Manager at Simula Research Laboratory \\
2001 -- 2005 & Scientific Programmer at Simula Research Laboratory \\
2000 -- 2001 & Scientific Programmer at Department of Informatics, University of Oslo \\
1996 -- 1999 & Research Assistant at Department of Informatics, University of Oslo \\
1996         & Summer intern, SINTEF Applied Mathematics \\
1993 -- 1994 & Assistant teacher, Department of Mathematics, University of Oslo \\
\end{tabular}

\section*{Languages}
\begin{tabular}{LR}
Norwegian & Mother tongue \\
English & Fluent \\
%German & Basic \\
%French & Rudimentary \\
\end{tabular}

\section*{Personal Skills}
\begin{tabular}{LR}
Management & Development and implementation of strategies, guidelines and
procedures, Lead and encourage staff, Personnel follow-up \\ \addlinespace
Communication & Writing and Editing, typesetting \& typography. Try to learn
French.
\\ \addlinespace
Funding & EU and RCN funding instruments, consortium development, proposal
writing \\ \addlinespace
Technology & Extensive knowledge about
system administration; always eager to learn about new frameworks and open
source technologies. Solution architect and problem solver. \\
\end{tabular}

\section*{Some Interests and Hobbies}
\begin{tabular}{LR}
Soft & Playing clarinet and saxophone, Cooking \& Wine, Photography \\
Hard & Bicycle, motorcycle \\
Tech & Fonts, typesetting with \TeX, gadgets, live coding 
\end{tabular}

\section*{Extended Description of Select Projects}

\begin{tabular}{@{}LR}
Activity & Tool for demographical and economical planning \\
Role & Architect and main developer \\
Staffing & Project leader, developer, front-end designer \\
Description & For the customer we developed a web-based system for statistical
estimation of demographics within a large closed group, and a tool for economical
planning of a large scale development. My main contribution was the mathematical
analysis of the problem and development of the necessary logic. The work was
carried out in close collaboration with financial experts and a designer.\\
Tools & Python, Pandas, Web.py, sqlite, MS SQL, Excel, HTML, CSS, JavaScript, jQuery \\
\addlinespace \bottomrule[.1pt] \addlinespace
\end{tabular}

\begin{tabular}{@{}LR}
Activity & Establishing technical infrastructure of Simula Research Laboratory \\
Role & Project leader and main architect \\
Staffing & Project leader, four system administrators, entrepreneur
representatives \\
Description & For the newly created research institute Simula Research
Laboratory we designed and developed a full technical infrastructure in the
renovated office building at the old airport at Fornebu; office network, backend
facilities and computer room, telephone systems, printers, servers. \\
\addlinespace \bottomrule[.1pt] \addlinespace
\end{tabular}

\begin{tabular}{@{}LR}
Activity & Symphonical \\
Role & Main developer \\
Staffing & Project leader and two developers \\
Description & Symphonical is a web-based collaboration tools. Initially it was
conceived as a tool for running agile software development projects based on the
metaphor of post-it notes on a virtual wall. We created the first prototype 
of the system, in 2005, before it was spun out as a company and devloped further
by others. \\
Tools & PHP, Mysql \\
\addlinespace \bottomrule[.1pt] \addlinespace
\end{tabular}

\begin{tabular}{@{}LR}
Activity & Simulation of Medical Ultrasound \\
Role & Researcher and main developer \\
Staffing & Four--five researchers and developers \\
Description & A research project where we studied different mathematical models
for numerical simulation of medical ultrasound. My main responsibility was
implementation of Finite element based simulators of ultrasound propagating in
inhomogeneous media. The target application was detection of tumors in tissue. \\
Tools & C++, Diffpack, MPI, VTK \\
\addlinespace \bottomrule[.1pt] \addlinespace
\end{tabular}

\begin{tabular}{@{}LR}
Activity & PySE \\
Role & Main developer \\
Description & Python Stencil Environment (PySE) is a high level environment for
solving partial differential equations using Finite difference methods. The
environment provides an abstraction for the stencil concept of FDM, which makes
it possible to implement FDM-based simulators in a notation very close to the
mathematical formulation. The environment handles parallelization and will
distribute the problem on the available hardware in an optimal way relatively to
the communication patterns defined by the stencils. PySE was
developed as part of my PhD studies and published as open source software.
\\
Tools & Python, Numpy, Scipy, SWIG, C++, MPI \\
\addlinespace \bottomrule[.1pt] \addlinespace
\end{tabular}

\begin{tabular}{@{}LR}
Activity & Python Computing Components (PyCC) \\
Role & Developer \\
Description & PyCC is a modern and efficient scripting framework that is used
to solve differential equations to study the electrical activity in the human
heart. The main parts of PyCC is a high level scripting interface implemented in
Python for flexibility and low-level extension modules implemented in C/C++ and
Fortran for efficiency. My main role in the project was to provide a flexible
build system for PyCC taking into account all the dependencies of the system.
The build system was implemented using the SCons framework. I was also involved
in \\
Tools & C/C++, Fortran, Python, Swig, MPI, Subversion, SCons, PETSc, FEniCS,
Hypre, BoomerAMG, Diffpack, GNU Compiler Collection and Debugger \\
\addlinespace \bottomrule[.1pt] \addlinespace
\end{tabular}

\begin{tabular}{@{}LR}
Activity & Web site for Simula Research Laboratory \\
Role & Main developer and designer \\
Description & In order to facilitates new communication strategies for the
company, it was necessary to update the web site for Simula Research Laboratory.
New channels for communicating news and ongoing research projects were created,
as well as a complete redesign of the company Intranet which at the same time
replaced mailing list as the primary form of official internal communication.
The process of upgrading the pages included both technical work and
meetings with all the stakeholders within the organisation.
\\
Tools & Python, PHP, Javascript, Plone, mysql \\
\addlinespace \bottomrule[.1pt] \addlinespace
\end{tabular}

\begin{tabular}{@{}LR}
Activity & Outsourcing of ICT services \\
Role & Project coordinator \\
Description & Simula Research Laboratory decided to put a contract for
its IT operations out for tender, after running the services with internal
resources for a decade. My core responsibility was authoring the specification
of the tender and evaluating the offers as well as take part in negotiating the
final contract. I also was responsible for handling the communication with all
the internal stakeholders. As the change affected the daily work of all
employees it was particularly important to have an open and inclusive dialog. \\
\addlinespace \bottomrule[.1pt] \addlinespace
\end{tabular}

\begin{tabular}{@{}LR}
Activity & \\
Role & \\
Staffing & \\
Description &  \\
Tools & \\
\addlinespace \bottomrule[.1pt] \addlinespace
\end{tabular}


Clustere

Google apps

IT Outsourcing

%VAMPiLE
%RITE



%\section*{Primary skills}
%
%Managing and Developing research funding proposals, for the framework programs of the
%European Commision, the European Research Council, and other research
%funding agencies.
%
%Personnel Management.
%
%Developing and managing administrative procedures.
%   
%Extensive technological experience, including web, social media, and
%backend/server programming. From my earlier PhD training, I have extensive
%knowledge about parallel programming and development of parallel computing
%systems. 
%
%Scientific scripting with Python and other tools.
%
%%Combining Python and C/C++ libraries using tools like Swig or SIP, or by
%%manually creating wrapper code.
%
%%Parallel programming with Python, using tools like pypar, mpi4py, pympi or
%%direct use of MPI through custom built C extensions.
%
%Solving practical system administration related problems on the Linux and
%MacOSX platform. 
%
%Based on my education and later involvement in the research at Simula Research
%Laboratory, I have a strong mathematical background. In my management work and
%proposal writing work in recent years, I have utilised this background to
%analyse and drill down into complex issues in several other fields, spanning
%from Scientific problems in related topics to development of budgeting
%procedures and administrative challenges. Although I am not working with
%mathematics every day, I am confident that my professional training has served
%me well and help me solve the various tasks I meet on a daily basis.
%
% \end{itemize}
% \item General Linux system administration. Have been a linux user since 1995, and a 
% professional linux administrator since 1999. I have mainly worked with Debian Linux,
% although I on several occasions also have looked into Gentoo, Suse and RedHat. 
% \item In recent years, I have also used and done system administration on the
% Mac OSX platform, and have gained significant insight on that platform as
% well. 
% \item Configuration and installation of Linux clusters. I have designed and
%installed three smaller linux clusters (from 32 to 96 processors), both
%hardware and software wise, based on Debian, FAI automatic installations,
%OpenPBS, Torque, and MAUI schedulers. 


\end{document}

% vim: tw=80
